\section{Conclusion}
\label{sec:conclusion}

We set out to build an identity and access platform that separates authentication, profile, tenant membership and authorization, and to prove with tests that its boundaries hold. The final revision does that. Keycloak proves who a person or program is; our backend decides, per tenant and per request, from facts read at decision time, what that principal may do. All seven milestones are implemented and verified, with M5 delivered as the agreed minimal provisioning slice. Both agent extensions are implemented: AG01 gives a program its own narrow, auditable identity, and AG02, finished in the final week as promised in the Week~1 design, lets a program act for a user with never more than that user's current permission, inside one tenant, until the user withdraws it.

The final clean-clone run of \texttt{e13f019} passed all 253 tests with real PostgreSQL and the live provider, killed all 15 deliberately injected security bugs, reported a stopped database as inconclusive instead of success, and passed the real-browser proof on two host names, including a delegation created by clicking a button in the browser and then used by the real agent.

The main lesson of the project is that the hard part of identity is not login, which a mature provider does well, but everything after it: deciding what a verified principal may do in a specific tenant right now, and proving it. The second lesson came from our mistakes: the evidence needs as much engineering as the code. Twice our tests could not fail, once a refactor removed a documented exception without any test noticing, and several times we reported results that we had not just observed. What fixed this was a discipline rather than a tool: every claim comes from a clean clone of a named commit, every guard has a mutation that must be killed, and every mistake is written down.
